\chapter{Proposed Multimodal Sensing-Assisted Beam Management Framework}
\label{ch:methodology}

This chapter presents the proposed BS-centric multimodal sensing-assisted beam
management framework. The objective of the framework is to use distributed
sensing observations, initial communication measurements, and temporal tracking
to support vehicle detection, per-target beam prediction, target-user
association, and beam maintenance in vehicular networks. The chapter focuses on
the algorithmic data flow and system interface. It does not claim measured
performance improvement, because the experimental evidence and the final
simulation settings are not yet fixed.

\section{Framework Overview}
\label{sec:method-framework}

For each serving base station (BS), the proposed framework constructs a local
cooperative sensing region consisting of the BS itself and \(K_{\mathrm{SN}}\) neighbouring
sensing nodes (SNs). At decision epoch \(t\), the BS receives multiview camera
observations from its own cameras and the selected SN cameras, together with a
BS-side ISAC point cloud. During initial access or explicit calibration periods,
the BS also obtains communication-side beam measurements, such as an SSB beam
sweep power spectrum. The framework converts these heterogeneous observations
into a unified bird's-eye-view (BEV) representation and then performs four
coupled tasks:
\begin{enumerate}
  \item vehicle target detection in the BEV coordinate system;
  \item per-target beam power distribution prediction and Top-\(K\) candidate
        beam selection;
  \item initial target-user association using SSB-spectrum information; and
  \item tracking-assisted beam maintenance using an interacting multiple model
        (IMM) filter.
\end{enumerate}

Let \(\mathcal{C}_{b,t}\) denote the set of camera images available to BS
\(b\) at epoch \(t\), including the BS cameras and the cameras of the selected
SNs. Let \(\mathcal{P}_{b,t}\) denote the BS-side ISAC point cloud, and let
\(\mathcal{Y}_{b,t}\) denote the available communication measurements. The
overall mapping can be written abstractly as
\begin{equation}
  \left(
  \mathcal{D}_{b,t},
  \widehat{\mathcal{K}}_{b,t},
  \mathcal{A}_{b,t},
  \mathcal{X}_{b,t}
  \right)
  =
  \mathcal{G}_{\Theta}
  \left(
  \mathcal{C}_{b,t},
  \mathcal{P}_{b,t},
  \mathcal{Y}_{b,t},
  \mathcal{X}_{b,t-1}
  \right),
  \label{eq:method-overall-mapping}
\end{equation}
where \(\mathcal{D}_{b,t}\) is the set of detected BEV targets,
\(\widehat{\mathcal{K}}_{b,t}\) is the set of predicted Top-\(K\) beam
candidates for the detected or tracked targets, \(\mathcal{A}_{b,t}\) is the
target-user association state, and \(\mathcal{X}_{b,t}\) is the set of tracking
states. The parameter set \(\Theta\) represents the sensing network, detection
head, beam prediction head, and tracking-related parameters. A final framework
figure should show this data flow from multimodal input construction to the
communication evaluation interface.

The framework distinguishes between sensing target identity and communication
user identity. A detected vehicle in the BEV space is first a sensing target.
It becomes a communication user only after the system establishes a mapping
between the detection or track and a user identity observed by communication
measurements. This distinction is important in multi-vehicle scenes, where the
strongest predicted beam for a target does not by itself identify the user that
generated a communication-side beam spectrum.

\section{BEV-Based Multimodal Fusion}
\label{sec:method-multimodal-processing}

The sensing front end maps multiview camera observations and the BS-side ISAC
point cloud into a common BEV coordinate frame. The purpose of this BEV
representation is to align observations from different viewpoints, station
locations, and modalities before target detection and beam prediction are
performed. For BS \(b\), the camera feature extractor first transforms each
image \(I_{v,t} \in \mathcal{C}_{b,t}\) into an image-plane feature map. These
features are then projected or lifted into BEV coordinates using the available
calibration information:
\begin{equation}
  \mathbf{F}^{\mathrm{cam}}_{b,t}
  =
  \Phi_{\mathrm{cam}}
  \left(
  \{ I_{v,t}, \mathbf{T}_{v \rightarrow b} \}_{v \in \mathcal{V}_{b}}
  \right),
  \label{eq:method-camera-bev}
\end{equation}
where \(\mathcal{V}_{b}\) is the set of camera views used by BS \(b\), and
\(\mathbf{T}_{v \rightarrow b}\) denotes the calibration transform from camera
view \(v\) to the BS-centered BEV frame.

The ISAC point cloud is encoded into the same BEV frame through a point-cloud
encoder:
\begin{equation}
  \mathbf{F}^{\mathrm{isac}}_{b,t}
  =
  \Phi_{\mathrm{isac}}
  \left(
  \mathcal{P}_{b,t};
  \mathbf{T}_{\mathrm{isac} \rightarrow b}
  \right).
  \label{eq:method-isac-bev}
\end{equation}
The camera and ISAC BEV features are subsequently fused into a joint BEV feature
map,
\begin{equation}
  \mathbf{F}_{b,t}
  =
  \Phi_{\mathrm{fuse}}
  \left(
  \mathbf{F}^{\mathrm{cam}}_{b,t},
  \mathbf{F}^{\mathrm{isac}}_{b,t},
  \mathbf{M}_{b,t}
  \right),
  \label{eq:method-bev-fusion}
\end{equation}
where \(\mathbf{M}_{b,t}\) denotes optional modality-availability or confidence
masks. Such masks are useful because distributed sensing observations may be
missing, delayed, partially occluded, or noisy.

The current design assumes that temporal synchronization and spatial
calibration are handled before BEV fusion. In practice, calibration and
synchronization errors can affect both detection and beam prediction. Therefore,
the implementation should keep explicit records of sensor timestamps, camera
extrinsics, point-cloud coordinate transforms, and missing-modality indicators.
TODO[unstable-method]: Finalize the BEV grid resolution, coordinate range,
camera-to-BEV projection method, ISAC point-cloud encoding, feature fusion
strategy, and the treatment of missing or noisy sensing observations.

\section{Vehicle Detection and Per-Target Beam Prediction}
\label{sec:method-detection-beam}

The joint BEV feature map \(\mathbf{F}_{b,t}\) is processed by a CenterNet-like
detection head. Instead of proposing a large set of anchor boxes, the detection
head predicts a BEV heatmap whose local maxima correspond to vehicle centers.
For each BEV grid cell, the head may also predict center offsets, target size,
orientation-related attributes, and a detection confidence score. The detection
output is written as
\begin{equation}
  \mathcal{D}_{b,t}
  =
  \left\{
  d_{i,t}
  =
  \left(
  \widehat{\mathbf{p}}_{i,t},
  \widehat{\mathbf{s}}_{i,t},
  q_{i,t},
  \mathbf{e}_{i,t}
  \right)
  \right\}_{i=1}^{N_{b,t}},
  \label{eq:method-detection-output}
\end{equation}
where \(\widehat{\mathbf{p}}_{i,t}\) is the BEV center position,
\(\widehat{\mathbf{s}}_{i,t}\) denotes target attributes such as size, \(q_{i,t}\)
is the detection confidence, \(\mathbf{e}_{i,t}\) is the target-level feature
embedding, and \(N_{b,t}\) is the number of detected targets.

For each detected target, the beam prediction head estimates a distribution over
the BS beam codebook
\(\mathcal{F}_{b}=\{\mathbf{f}_{b,1},\ldots,\mathbf{f}_{b,M_b^{\mathrm{beam}}}\}\). The output
for target \(i\) is
\begin{equation}
  \widehat{\mathbf{g}}_{i,t}
  =
  \left[
  \widehat{g}_{i,t,1},
  \widehat{g}_{i,t,2},
  \ldots,
  \widehat{g}_{i,t,M_b^{\mathrm{beam}}}
  \right],
  \label{eq:method-beam-distribution}
\end{equation}
where \(\widehat{g}_{i,t,m}\) represents the predicted quality, normalized
power, or score of beam \(m\) for the detected target. The Top-\(K\) candidate
beam set is then defined as
\begin{equation}
  \widehat{\mathcal{K}}_{i,t}^{(K)}
  =
  \underset{
  \mathcal{S} \subseteq \{1,\ldots,M_b^{\mathrm{beam}}\},\ |\mathcal{S}|=K_{\mathrm{beam}}
  }{\operatorname{arg\,max}}
  \sum_{m \in \mathcal{S}} \widehat{g}_{i,t,m}.
  \label{eq:method-topk-beams}
\end{equation}
This distributional output is more informative than a single hard beam index:
it supports Top-\(K\) beam selection, uncertainty-aware association, and
communication-side mismatch modelling.

The training objective is expected to combine detection and beam-prediction
terms,
\begin{equation}
  \mathcal{L}
  =
  \lambda_{\mathrm{det}}\mathcal{L}_{\mathrm{det}}
  +
  \lambda_{\mathrm{beam}}\mathcal{L}_{\mathrm{beam}}
  +
  \lambda_{\mathrm{aux}}\mathcal{L}_{\mathrm{aux}},
  \label{eq:method-training-loss}
\end{equation}
where \(\mathcal{L}_{\mathrm{det}}\) may include heatmap, offset, and attribute
losses, \(\mathcal{L}_{\mathrm{beam}}\) supervises the predicted per-beam
distribution, and \(\mathcal{L}_{\mathrm{aux}}\) denotes optional auxiliary
terms. TODO[unstable-method]: Finalize the exact detection loss, beam
distribution loss, target-to-label assignment, positive and negative sample
definition, beam-label generation process, and beam power normalization.

\section{Initial Target-User Association}
\label{sec:method-association}

Initial target-user association is performed when communication-side beam
measurements are available. During initial access, the BS conducts an SSB beam
sweep and obtains a beam power spectrum for each communication user. Let
\(\mathbf{y}_{u,t_0}\) denote the observed spectrum for user \(u\) at the
initial association epoch \(t_0\):
\begin{equation}
  \mathbf{y}_{u,t_0}
  =
  \left[
  y_{u,t_0,1},
  y_{u,t_0,2},
  \ldots,
  y_{u,t_0,M_b^{\mathrm{beam}}}
  \right].
  \label{eq:method-ssb-spectrum}
\end{equation}
The sensing network simultaneously predicts
\(\widehat{\mathbf{g}}_{i,t_0}\) for each detected target \(i\). Association is
then formulated as a matching problem between detected sensing targets and
communication user identities.

A generic association cost can be expressed as
\begin{equation}
  c_{i,u}
  =
  \alpha_{\mathrm{beam}}
  d_{\mathrm{beam}}
  \left(
  \widehat{\mathbf{g}}_{i,t_0},
  \widetilde{\mathbf{y}}_{u,t_0}
  \right)
  +
  \alpha_{\mathrm{geo}}
  d_{\mathrm{geo}}(i,u)
  -
  \alpha_{\mathrm{conf}} q_{i,t_0},
  \label{eq:method-association-cost}
\end{equation}
where \(\widetilde{\mathbf{y}}_{u,t_0}\) is a normalized SSB spectrum,
\(d_{\mathrm{beam}}\) measures distributional disagreement,
\(d_{\mathrm{geo}}\) represents any available geometric prior, and \(q_{i,t_0}\)
is the detection confidence. The association variable is
\(a_{i,u,t_0} \in \{0,1\}\), where \(a_{i,u,t_0}=1\) means that sensing target
\(i\) is assigned to communication user \(u\). A one-to-one assignment can be
written as
\begin{equation}
  \begin{aligned}
  \min_{\{a_{i,u,t_0}\}} \quad
  & \sum_{i=1}^{N_{b,t_0}}\sum_{u=1}^{U_{b,t_0}}
    a_{i,u,t_0} c_{i,u} \\
  \mathrm{s.t.}\quad
  & \sum_{u=1}^{U_{b,t_0}} a_{i,u,t_0} \leq 1,\quad
    \sum_{i=1}^{N_{b,t_0}} a_{i,u,t_0} \leq 1, \\
  & a_{i,u,t_0} \in \{0,1\}.
  \end{aligned}
  \label{eq:method-initial-assignment}
\end{equation}
This formulation gives the intended interface without fixing the final cost
terms. TODO[unstable-method]: Finalize the association cost, spectrum
normalization, geometric prior, matching algorithm, conflict handling, rejection
thresholds, and recovery behaviour when no reliable one-to-one match exists.

Once a match is accepted, the framework creates a unified track containing a
sensing track identity, a communication user identity, a BEV state estimate, and
the current Top-\(K\) beam candidates. These associations are then maintained by
the tracking module rather than by repeated full beam sweeps at every epoch.

\section{IMM Tracking and Beam Maintenance}
\label{sec:method-tracking}

After initial association, the framework maintains each associated user through
an IMM tracking module. The tracking state for track \(\ell\) at epoch \(t\) is
represented as
\begin{equation}
  \mathbf{x}_{\ell,t}
  =
  \left[
  x_{\ell,t},
  y_{\ell,t},
  \dot{x}_{\ell,t},
  \dot{y}_{\ell,t}
  \right]^{\mathrm{T}},
  \label{eq:method-track-state}
\end{equation}
where \((x_{\ell,t},y_{\ell,t})\) is the BEV position and
\((\dot{x}_{\ell,t},\dot{y}_{\ell,t})\) is the BEV velocity. Additional
components, such as acceleration, heading, or lane-related state variables, can
be introduced if the final simulator and tracking design require them.

The IMM filter maintains a set of candidate motion models
\(\mathcal{Q}=\{1,\ldots,Q\}\). For each mode \(q \in \mathcal{Q}\), the predicted state is
abstractly written as
\begin{equation}
  \mathbf{x}^{(q)}_{\ell,t|t-1}
  =
  \mathbf{F}^{(q)}
  \mathbf{x}^{(q)}_{\ell,t-1|t-1},
  \qquad
  \mathbf{P}^{(q)}_{\ell,t|t-1}
  =
  \mathbf{F}^{(q)}
  \mathbf{P}^{(q)}_{\ell,t-1|t-1}
  \left(\mathbf{F}^{(q)}\right)^{\mathrm{T}}
  +
  \mathbf{Q}^{(q)}.
  \label{eq:method-imm-prediction}
\end{equation}
The mode-conditioned predictions are mixed according to the mode probabilities,
and the resulting predicted track is associated with the next-frame BEV
detections. If detection \(i\) provides measurement
\(\mathbf{z}_{i,t}=\widehat{\mathbf{p}}_{i,t}\), a generic tracking-association
score can include innovation distance, motion consistency, detection confidence,
and beam-distribution consistency:
\begin{equation}
  g_{\ell,i,t}
  =
  \beta_{\mathrm{pos}}
  d_{\mathrm{pos}}
  \left(
  \mathbf{z}_{i,t},
  \widehat{\mathbf{x}}_{\ell,t|t-1}
  \right)
  +
  \beta_{\mathrm{beam}}
  d_{\mathrm{beam}}
  \left(
  \widehat{\mathbf{g}}_{i,t},
  \widehat{\mathbf{g}}_{\ell,t-1}
  \right)
  -
  \beta_{\mathrm{conf}} q_{i,t}.
  \label{eq:method-track-detection-cost}
\end{equation}
The accepted detection updates the track state, refreshes the predicted beam
distribution, and propagates the communication user identity. If no reliable
detection is matched, the framework may keep a prediction-only track for a
limited duration and trigger a partial beam measurement or re-association
procedure.

Beam maintenance uses both the current beam prediction and the temporal
continuity encoded by the track. For a matched track \(\ell\), the communication
module receives the candidate set
\(\widehat{\mathcal{K}}_{\ell,t}^{(K)}\) from the associated detection. For a
temporarily unmatched track, the framework may use the predicted state and the
previous beam distribution as a short-term fallback, subject to confidence and
age limits. TODO[unstable-method]: Finalize the IMM motion-model set, mode
transition matrix, process and observation noise parameters, gating thresholds,
track birth and deletion rules, prediction-only lifetime, and the exact trigger
for partial beam measurement or re-association.

\section{Communication Interface and System-Level Evaluation}
\label{sec:method-communication-interface}

The communication interface receives, for each associated user \(u\), a user
identity, a tracking state, and a Top-\(K\) candidate beam set. These outputs can
be used in two ways. First, they can reduce the candidate set over which CSI
acquisition or PMI selection is performed. Second, they provide a structured
input to the system-level evaluation model that estimates effective SINR,
spectral efficiency, and throughput under sensing-assisted beam management.

Let \(m_{u,t}^{\star}\) denote the best beam index under the generated channel
or beam-label model, and let
\(\widehat{\mathcal{K}}_{u,t}^{(K)}\) denote the Top-\(K\) candidate set
provided by the sensing and tracking pipeline. The coverage indicator is
\begin{equation}
  \delta_{u,t}^{(K)}
  =
  \begin{cases}
  1, & m_{u,t}^{\star} \in \widehat{\mathcal{K}}_{u,t}^{(K)},\\
  0, & \text{otherwise}.
  \end{cases}
  \label{eq:method-topk-coverage}
\end{equation}
If \(\delta_{u,t}^{(K)}=1\), the evaluation can assume that subsequent
communication measurements are concentrated on the candidate set rather than on
the full codebook. If \(\delta_{u,t}^{(K)}=0\), the evaluation introduces a
beam-mismatch penalty to model imperfect CSI or PMI selection.

For a multi-user ZF abstraction, let \(\widetilde{\mathbf{h}}_{u,t}\) be the
effective channel used for precoding after the candidate-beam decision. The
stacked channel matrix is
\begin{equation}
  \widetilde{\mathbf{H}}_{t}
  =
  \left[
  \widetilde{\mathbf{h}}_{1,t}^{\mathrm{T}};
  \widetilde{\mathbf{h}}_{2,t}^{\mathrm{T}};
  \ldots;
  \widetilde{\mathbf{h}}_{U_t,t}^{\mathrm{T}}
  \right],
  \label{eq:method-effective-channel}
\end{equation}
and the ZF precoder can be represented abstractly by the pseudo-inverse of
\(\widetilde{\mathbf{H}}_{t}\), followed by the chosen power normalization. The
effective SINR for user \(u\) can then be expressed as
\begin{equation}
  \gamma_{u,t}
  =
  \frac{
  P_u
  \left|
  \mathbf{h}_{u,t}^{\mathrm{H}}\mathbf{w}_{u,t}
  \right|^2
  }{
  \sum_{v \neq u}
  P_v
  \left|
  \mathbf{h}_{u,t}^{\mathrm{H}}\mathbf{w}_{v,t}
  \right|^2
  +
  \sigma^2
  },
  \label{eq:method-zf-sinr}
\end{equation}
where \(\mathbf{w}_{u,t}\) is the precoding vector assigned to user \(u\),
\(P_u\) is the allocated transmit power, and \(\sigma^2\) is the noise power.
The spectral efficiency is
\begin{equation}
  R_{u,t}
  =
  \log_2(1+\gamma_{u,t}),
  \label{eq:method-spectral-efficiency}
\end{equation}
and an overhead-adjusted sum throughput can be written as
\begin{equation}
  R^{\mathrm{adj}}_{t}
  =
  (1-\rho_t)
  \sum_{u=1}^{U_t} R_{u,t},
  \label{eq:method-overhead-adjusted-throughput}
\end{equation}
where \(\rho_t\) denotes the fraction of time-frequency resources consumed by
beam training, calibration, or recovery measurements at epoch \(t\).

These equations define an evaluation interface rather than a finalized
communication protocol. TODO[unstable-method]: Finalize the beam mismatch
penalty, how mismatch modifies CSI or PMI abstraction, ZF normalization,
overhead accounting, and the exact mapping from Top-\(K\) coverage to SINR
degradation.

\section{Algorithm Summary and Complexity}
\label{sec:method-algorithm-complexity}

Algorithm~\ref{alg:method-framework} summarizes the online operation of the
proposed framework. The procedure separates the initial access stage, where SSB
spectra are used to bind sensing targets to communication users, from the
subsequent maintenance stage, where IMM tracking and BEV detections preserve
identity continuity and reduce the need for full beam sweeps.

\begin{algorithm}[t]
\caption{Multimodal sensing-assisted beam management at BS \(b\)}
\label{alg:method-framework}
\begin{algorithmic}[1]
\REQUIRE Multiview camera inputs \(\mathcal{C}_{b,t}\), ISAC point cloud
\(\mathcal{P}_{b,t}\), beam measurements \(\mathcal{Y}_{b,t}\) when available,
previous tracks \(\mathcal{X}_{b,t-1}\), beam codebook \(\mathcal{F}_{b}\)
\ENSURE Detections \(\mathcal{D}_{b,t}\), Top-\(K\) beam sets
\(\widehat{\mathcal{K}}_{b,t}\), associations \(\mathcal{A}_{b,t}\), updated
tracks \(\mathcal{X}_{b,t}\)
\STATE Construct camera BEV features using calibrated BS/SN camera views.
\STATE Encode the BS-side ISAC point cloud in the same BEV coordinate frame.
\STATE Fuse camera and ISAC BEV features to obtain \(\mathbf{F}_{b,t}\).
\STATE Detect vehicle targets and extract target-level features from
\(\mathbf{F}_{b,t}\).
\STATE Predict a per-beam distribution and Top-\(K\) candidate beams for each
detected target.
\IF{epoch \(t\) is an initial access or re-association epoch}
  \STATE Obtain SSB beam power spectra from communication users.
  \STATE Match detected targets to communication user identities using the
  association cost.
  \STATE Initialize or repair tracks with matched target-user pairs.
\ELSE
  \STATE Predict existing tracks using the IMM filter.
  \STATE Associate new BEV detections with predicted tracks.
  \STATE Update matched tracks and refresh their Top-\(K\) beam sets.
  \STATE Mark unmatched tracks for prediction-only maintenance, partial beam
  measurement, or re-association according to confidence rules.
\ENDIF
\STATE Pass user identities, track states, and Top-\(K\) beam sets to the
communication evaluation interface.
\STATE Compute beam-training overhead, effective SINR, spectral efficiency, and
overhead-adjusted throughput when channel labels are available.
\end{algorithmic}
\end{algorithm}

The main online cost consists of five parts. First, BEV feature construction and
fusion dominate the sensing front end and depend on the number of camera views,
the point-cloud size, and the BEV grid resolution. Second, the detection and
beam heads add target-level inference cost; the beam head scales with the number
of detected targets and the codebook size \(M_b^{\mathrm{beam}}\). Third, initial target-user
association requires matching \(N_{b,t}\) detected targets to \(U_{b,t}\)
communication users. If a standard bipartite assignment solver is used after
the association cost is fixed, this step scales cubically in the larger side of
the matching problem. Fourth, IMM tracking scales with the number of active
tracks and the number of motion modes. Finally, the communication evaluation
cost depends on the number of associated users and the matrix operations used by
the ZF abstraction.

This complexity discussion is intentionally qualitative at the current stage.
TODO[needs-evidence]: Add measured runtime, memory usage, and implementation
details after the model architecture, simulator, and evaluation scripts are
fixed.
